\section{Related Work}
\noindent{\textbf{End-to-End Autonomous Driving.}}
Conventional autonomous systems decompose the whole paradigm into three stages:  perception~\cite{li2024bevformer,wang2024not,wang2023exploring,liu2024mgmap,liu2025uncertainty}, prediction~\cite{song2020pip,xin2018intention,liu2021survey}, and planning~\cite{lim2019hybrid,gonzalez2015review}. Such modular designs often lead to error accumulation and substantial engineering overhead. 
To address these limitations, recent works explore end-to-end autonomous driving, which learns a direct mapping from sensory observations to driving actions or trajectories~\cite{hu2023planning,sun2025sparsedrive,li2024ego,jiang2023vad}. Existing methods can be broadly categorized 
based on the form of action outputs, including regression-based methods~\cite{hu2022st, jiang2023vad}, generation-based approaches~\cite{zheng2024genad,liao2025diffusiondrive}, and reinforcement-based methods~\cite{gao2025rad,yan2025adr1}. In regression-based approaches, UniAD~\cite{hu2023planning}, VAD~\cite{jiang2023vad}, and SparseDrive~\cite{sun2025sparsedrive} construct end-to-end architectures by leveraging BEV representations and vectorized or sparse queries. For generative methods, GenAD~\cite{zheng2024genad} leverages generative models for trajectory regression. DiffusionDrive~\cite{liao2025diffusiondrive} and GoalFlow~\cite{xing2025goalflow} utilize flow matching~\cite{lipman2022flow} and diffusion policy~\cite{chi2025diffusion} paradigms to sample multi-mode actions. For reinforcement learning-based approaches, recent studies~\cite{li2025driver1,jiao2025evadrive, li2025recogdrive} formulate autonomous driving as a sequential decision-making problem, which optimizes driving policies through reward design and interaction with the environment. In addition, some recent approaches incorporate large language models (LLMs) into vision-language-action (VLA) frameworks by injecting language context as high-level guidance for reasoning and decision-making~\cite{li2025recogdrive,chi2025impromptu,zhou2025opendrivevla}. 

However, most existing end-to-end driving methods lack the ability to reason about future scene dynamics under different driving actions. 
This limitation motivates the development of world models that simulate future world evolution to support more informed decision-making.

\noindent{\textbf{World Models for Autonomous Driving}}
In autonomous driving systems, the world model aims to predict the future evolution of scenes following various actions~\cite{guan2024world,kong20253d,wang2024drivingfuture,zheng2025world4drive,li2024law}. 
Combining with the next state prediction, existing world models can be categorized into two types: future scene prediction~\cite{wang2024drivingfuture,zhang2025epona,zhao2025drivedreamer4d,wang2025uniocc,zheng2024occworld} and feature-level world models in latent space~\cite{li2024law,li2025wote,li2024ssr,yang2025worldrft}. 
Representative works of future scene prediction focus on directly forecasting future observations or structured scene representations. 
For instance, GenAD~\cite{yang2024generalized}, DrivingFuture~\cite{wang2024drivingfuture}, and Epona~\cite{zhang2025epona} predict future pixel-level scenes conditioned on actions, while DriveDreamer4D~\cite{zhao2025drivedreamer4d} extends this paradigm to 4D dynamic scene generation with 4D Gaussian Splatting reconstruction. 
Other approaches like UniOcc~\cite{wang2025uniocc} and OccWorld~\cite{zheng2024occworld} operate in voxelized or occupancy space, predicting future 3D occupancy or semantic fields to model scene evolution.
Complementary to explicit scene forecasting, another line of research efforts learns the world model in a compact latent space, where the model predicts action-conditioned transitions over intermediate features rather than reconstructing full observations. 
LAW~\cite{li2024law} and World4Drive~\cite{zheng2025world4drive} construct the world model in latent feature space for planning-oriented prediction, while SSR~\cite{li2024ssr} and WoTE~\cite{li2025wote} further explore latent modeling in BEV space, moving beyond perspective-view representations.

For our design, DynFlowDrive avoids explicit pixel-level scene generation and instead operates in latent space, focusing on modeling the underlying dynamics of world evolution for action-conditioned prediction.




